\documentclass[11pt,a4paper]{report}

% ----------------------------------------------------------------------
% Project Handover Document
% UATMS — Ultra Advanced Threat Monitoring System (Big Chemistry XDR)
% ----------------------------------------------------------------------

\usepackage[utf8]{inputenc}
\usepackage[T1]{fontenc}
\usepackage{lmodern}
\usepackage{geometry}
\geometry{a4paper, margin=2.5cm}
\usepackage{graphicx}
\usepackage{float}
\usepackage{xcolor}
\usepackage{titlesec}
\usepackage{enumitem}
\usepackage{fancyhdr}
\usepackage{booktabs}
\usepackage{hyperref}

% --- Colours ---------------------------------------------------------
\definecolor{primary}{HTML}{1F3A5F}
\definecolor{accent}{HTML}{D9772B}
\definecolor{lightgray}{HTML}{F2F2F2}

\hypersetup{
    colorlinks=true,
    linkcolor=primary,
    urlcolor=accent,
    citecolor=primary,
    pdftitle={UATMS Project Handover Document},
    pdfauthor={Big Chemistry XDR Team}
}

% --- Chapter / section styling --------------------------------------
\titleformat{\chapter}[display]
  {\normalfont\huge\bfseries\color{primary}}
  {\filright\Large\color{accent}\MakeUppercase{\chaptertitlename}\ \thechapter}
  {1ex}
  {\titlerule\vspace{1ex}\filright}

\titleformat{\section}
  {\normalfont\Large\bfseries\color{primary}}
  {\thesection}{1em}{}

\titleformat{\subsection}
  {\normalfont\large\bfseries\color{accent}}
  {\thesubsection}{1em}{}

% --- Header / footer -------------------------------------------------
\pagestyle{fancy}
\fancyhf{}
\renewcommand{\headrulewidth}{0.4pt}
\renewcommand{\footrulewidth}{0.2pt}
\fancyhead[L]{\small\color{primary}UATMS — Project Handover}
\fancyhead[R]{\small\color{primary}Big Chemistry XDR}
\fancyfoot[C]{\thepage}

% ====================================================================
\begin{document}

% --- Title page ------------------------------------------------------
\begin{titlepage}
    \centering
    \vspace*{2cm}

    {\color{accent}\rule{\textwidth}{2pt}}\\[0.6cm]
    {\Huge\bfseries\color{primary} Project Handover Document\par}
    \vspace{0.4cm}
    {\LARGE\color{primary} RobotLab's Advanced Threat Monitoring System \par}
    \vspace{0.2cm}
    {\large Big Chemistry — Extended Detection \& Response (XDR) Platform\par}
    {\color{accent}\rule{\textwidth}{2pt}}\\[1.5cm]

    \begin{tabular}{rl}    
        \textbf{Date:}          & 12 June 2026 \\[4pt]
    \end{tabular}

    \vfill
    
    \vspace{1cm}
\end{titlepage}

% --- TOC -------------------------------------------------------------
\tableofcontents
\thispagestyle{fancy}

% ====================================================================
\chapter{Current State of the Project}

This document hands over the RobotLab's Advanced Threat Monitoring System,
a full Extended Detection \& Response (XDR) platform built for Big Chemistry on
AWS \texttt{eu-central-1}. It captures the system as it stands today and lays out
the recommended order of work for the team taking it forward.

\section{Architecture Overview}

The platform follows a two-VPC connected by VPC Peering,
with no Transit Gateway. The control plane (\textbf{bc-ctrl}, \texttt{10.0.0.0/16})
hosts EC2-only workloads — Wazuh SIEM, Splunk SOAR, MISP, and the self-hosted
GitHub runner. The production plane (\textbf{bc-prd}, \texttt{10.30.0.0/16}) runs
the EKS cluster carrying the security stack and sensor DaemonSets (Cilium, Falco,
Tetragon, Zeek, Suricata, plus the NOMAD Oasis observability target).

The diagram below reflects the current, deployed state of the infrastructure.

\begin{figure}[H]
    \centering
    \includegraphics[width=\textwidth]{assets/Architecture_Diagram.jpg}
    \caption{Current-state architecture of the XDR platform on AWS
    \texttt{eu-central-1}. bc-ctrl (the ``Brain'') and bc-prd (the ``Data'')
    VPCs are joined by VPC Peering; shared services (S3, IAM, Route53, Lambda)
    sit beneath both.}
    \label{fig:architecture}
\end{figure}

\section{What Is Live Today}

\begin{itemize}[leftmargin=*]
    \item \textbf{bc-ctrl (Control plane):} Wazuh all-in-one (Manager + Indexer +
          Dashboard), MISP threat-intel platform, Splunk SOAR, and
          a self-hosted GitHub Actions runner.
    \item \textbf{bc-prd (Production plane):} EKS cluster running Cilium , Falco (\texttt{modern\_ebpf}), Tetragon, Hubble,
          and the Zeek / Suricata / wazuh-agent DaemonSets.
    \item \textbf{Telemetry pipeline:} Suricata, Zeek, Falco, and Tetragon emit to
          host-path logs that the wazuh-agent ships to the Wazuh manager. 
    \item \textbf{Threat-intel sync:} MISP IOCs are synced hourly into Zeek Intel
          and Suricata rule formats via sidecar containers.
\end{itemize}

The platform is operational but remains in a pre-production hardening phase. The
chapters that follow define the recommended sequence of work.

% ====================================================================
\chapter{First Priority: Stabilise NOMAD Oasis}

\section{Why This Comes First}

The single most important first task for the incoming team is to ensure that
\textbf{NOMAD Oasis is running properly}. NOMAD Oasis is the primary observability
and workload target on the bc-prd EKS cluster — it is what exercises the security
stack end-to-end. Until it is reliably healthy, every downstream detection,
enforcement, and analytics effort is being validated against an unstable base.

To date, NOMAD Oasis has converged on the cluster but has \emph{not} reached a
fully validated, repeatable healthy state.

%\section{Concrete Actions}

%\begin{enumerate}[leftmargin=*]
%    \item \textbf{Confirm a clean cold-start.} Run a full pipeline deployment from
 %         empty AWS state and verify NOMAD Oasis reaches a green, ready state with
  %        zero manual intervention. Capture the run ID and pod-readiness output as
   
          %evidence.
    %\item \textbf{Validate the supporting data tier.} MongoDB, PostgreSQL,
    %      Elasticsearch, RabbitMQ, and Temporal must all be healthy and bound to
          %their PVCs (gp3 / EFS). Confirm Temporal connects to Postgres and the
          %schema job completes.
   % \item \textbf{Verify CiliumNetworkPolicy selectors.} Mixed Helm sub-charts emit
   %       inconsistent labels; confirm every CNP selects the pods it intends to
    %      (use \texttt{kubectl get pod -o jsonpath='\{.metadata.labels\}'} and match
     %     exactly). A wrong selector silently drops traffic under
          %\texttt{policyEnforcementMode=default}.
    %\item \textbf{Resolve the Keycloak browser-login blocker.} The dev-path HTTP
    %      login is blocked by Keycloak 24's \texttt{SameSite=None; Secure} cookie
     %     defaults. The recommended fix is the production path: ALB Ingress + ACM
      %    certificate + Route53, giving real HTTPS so secure cookies work natively.
    %\item \textbf{Reconcile node-count drift.} Documentation states 2× t3.medium;
     %     the live cluster currently runs additional nodes (workload pool plus the
      %    \texttt{dedicated=nomad} node). Confirm the node-group configuration
       %   matches reality and is intentional.
%\end{enumerate}

\noindent\fcolorbox{accent}{lightgray}{%
\begin{minipage}{0.97\textwidth}
\textbf{Exit criterion:} A documented, reproducible cold-start that brings NOMAD
Oasis to a fully healthy state (including browser login) with no manual steps.
Only once this is met should the team move on to the hardening work below.
\end{minipage}}

% ====================================================================
\chapter{Unit-Scope Security Hardening}

With NOMAD Oasis stable, the next focus shifts to a more \textbf{unitary scope}:
embedding security controls at the level of individual code and container
artifacts, rather than only at the runtime / network layer where the platform is
already strong. These controls close the supply-chain and application-security
gaps that the current detection-centric posture does not cover.

\section{Container Image Scanning}

The platform currently has no container image scanning or admission control
, so images are deployed without provenance or vulnerability
verification. Recommended measures:

\begin{itemize}[leftmargin=*]
    \item \textbf{Vulnerability scanning} of every image at build time and in the
          registry (e.g. Trivy or Grype against ECR), failing the pipeline on
          critical / high CVEs.
    \item \textbf{Image signing and verification} (e.g. Cosign / Sigstore) so only
          signed, attested images can run.
    \item \textbf{Admission control} on the cluster (e.g. a policy controller) to
          enforce signature and scan-result requirements before a pod is admitted.
\end{itemize}

\section{Static Application Security Testing (SAST)}

Introduce SAST into the CI pipeline to catch vulnerabilities in source before
they ship — insecure patterns, injection flaws, hardcoded secrets, and unsafe
dependencies. Run it on every pull request and gate merges on the results. This
complements the existing infrastructure review (\texttt{/tf-review},
\texttt{/security-audit}) with application-code coverage.

\section{Dynamic Application Security Testing (DAST)}

Add DAST against running services (the NOMAD Oasis web tier, the Wazuh / MISP
dashboards, Splunk SOAR, and any exposed API). DAST exercises the deployed
application from the outside to surface authentication weaknesses, injection,
and misconfiguration that only manifest at runtime.

\section{Fuzzing}

For the higher-risk parsing and ingestion surfaces — the telemetry pipeline that
consumes Suricata \texttt{eve.json}, Zeek logs, and MISP-sourced rules — introduce
\textbf{fuzz testing} to harden input handling against malformed or adversarial
data. Untrusted external data crossing a system boundary is exactly where
fuzzing yields the most value.

\noindent\fcolorbox{accent}{lightgray}{%
\begin{minipage}{0.97\textwidth}
\textbf{Goal:} shift security ``left'' to the individual code and container unit,
covering SAST, DAST, and fuzzing in CI plus image scanning and signing at the
registry / admission boundary — closing the supply-chain gaps the runtime stack
cannot.
\end{minipage}}

% ====================================================================
\chapter{Forward Research: A Machine Learning Pipeline for XDR}

The next strategic direction for the research effort is to build a
\textbf{Machine Learning pipeline for XDR}. The platform already aggregates a
rich, multi-source telemetry stream — Suricata, Zeek, Falco, Tetragon, Wazuh
alerts, VPC flow logs, CloudTrail, and GuardDuty — which is precisely the
high-volume, labelled-and-unlabelled data foundation an ML detection layer needs.

\section{Motivation}

Today the system is rule- and signature-driven: it detects what it has been
explicitly told to look for. An ML layer adds the ability to surface
\emph{unknown} threats — anomalous process trees, unusual network flows,
credential-misuse patterns, and slow-burn lateral movement — that evade static
rules. It also offers a path to reducing alert fatigue by scoring and correlating
the existing detections.

%\section{Proposed Pipeline Stages}

%\begin{enumerate}[leftmargin=*]

    %\item \textbf{Data collection \& lake.} Centralise the existing telemetry into
     %     a queryable store (the S3 log buckets are the natural landing zone) with
      %    a consistent, normalised schema across sensors.
    %\item \textbf{Feature engineering.} Derive behavioural features — per-host
     %     process baselines, network-flow statistics, authentication sequences,
      %    and time-series aggregates — from the normalised events.
    %\item \textbf{Model development.} Start with unsupervised anomaly detection
    %      (the realistic first step given limited labelled incidents), then layer
    %      supervised models as labelled alert outcomes accumulate from analyst
    %      triage.
    %\item \textbf{Inference \& integration.} Serve models against the live stream
    %      and feed scored detections back into Wazuh / Splunk SOAR so they enrich
    %      existing alerts and can drive automated response.
    %\item \textbf{Feedback loop \& MLOps.} Capture analyst verdicts as labels,
    %      monitor for model and data drift, and retrain on a schedule — closing
    %      the loop so the system improves over time.
%\end{enumerate}

\section{Considerations}

\begin{itemize}[leftmargin=*]
    \item \textbf{Cost discipline.} ML infrastructure (training compute,
          storage, inference) must be scoped and justified.
    \item \textbf{Explainability.} Detections that drive response must be
          explainable to stakeholders; favour models and scoring whose output a human
          can reason about during incident response.
    \item \textbf{Data governance.} Telemetry includes sensitive operational
          detail; apply the same least-privilege and KMS-encryption discipline used
          across the rest of the platform to the data lake and ML artifacts.
\end{itemize}

\noindent\fcolorbox{accent}{lightgray}{%
\begin{minipage}{0.97\textwidth}
\textbf{Direction:} evolve the monitoring system from a rule-driven detection platform into a
learning system — an ML pipeline that turns the existing multi-sensor telemetry
into adaptive, anomaly-aware XDR detection.
\end{minipage}}

% ====================================================================
\chapter{Summary Roadmap}

\begin{table}[H]
    \centering
    \renewcommand{\arraystretch}{1.4}
    \begin{tabular}{@{}p{1.2cm}p{4.5cm}p{8cm}@{}}
        \toprule
        \textbf{Order} & \textbf{Workstream} & \textbf{Outcome} \\
        \midrule
        1 & Stabilise NOMAD Oasis & Reproducible, fully healthy cold-start
            including browser login; stable base for all downstream work. \\
        2 & Unit-scope security hardening & Container scanning + signing +
            admission control, and SAST / DAST / fuzzing in CI. \\
        3 & ML pipeline for XDR & Adaptive, anomaly-aware detection built on the
            existing multi-sensor telemetry. \\
        \bottomrule
    \end{tabular}
    \caption{Recommended order of work for the incoming team.}
\end{table}

\vspace{1cm}
\begin{center}
    {\color{primary}\itshape End of Handover Document.}
\end{center}

\end{document}
